\documentclass[10pt,a4paper]{amsart}
\usepackage[margin=19mm]{geometry}
\usepackage{amsmath,amssymb,mathtools}
\usepackage[scr=rsfs,cal=euler]{mathalfa}
\usepackage{xfrac}
\usepackage[unicode,hidelinks,bookmarks=false]{hyperref}
\newcommand{\ie}{i.e.\ }
\newcommand{\cf}{cf.\ }
\newcommand{\resp}{respectively\ }
\newcommand{\Subsection}{\subsection}
\providecommand{\nobreakdash}{\nobreak\hbox{-}}
\setlength{\emergencystretch}{2em}
\multlinegap0pt
\usepackage{bm}
\usepackage{bbm}
\usepackage{dsfont}
\usepackage{xspace}
\usepackage{cancel}
\usepackage{mathtools}
\usepackage{amsthm}
\makeatletter
\@ifundefined{theorem}{\newtheorem{theorem}{Theorem}}{}
\@ifundefined{lemma}{\newtheorem{lemma}[theorem]{Lemma}}{}
\@ifundefined{proposition}{\newtheorem{proposition}[theorem]{Proposition}}{}
\makeatother
\DeclareMathOperator{\End}{\mathcal{E}nd}
\newcommand{\Aff}{\mathbb{A}}
\newcommand{\NN}{\mathbb{N}}
\newcommand{\OO}{\mathcal{O}}
\newcommand{\PP}{\mathbb{P}}
\newcommand{\cD}{\mathcal{D}}
\newcommand{\cH}{\mathcal{H}}
\newcommand{\cL}{\mathcal{L}}
\newcommand{\cV}{\mathcal{V}}
\newcommand{\op}{\operatorname}
\renewcommand{\tilde}{\widetilde}		%never use ordinary tilde
\title[Towards Brill-Noether Theory for Spectral Curves]{Towards Brill-Noether Theory for Spectral Curves}
\author{Clemens Nollau}
\date{Source: arXiv:2408.11404v2 (CC BY 4.0)}
\begin{document}
\maketitle

\noindent\emph{Source and licence.} Towards Brill-Noether Theory for Spectral Curves. Version arXiv:2408.11404v2 (CC BY 4.0), distributed under the Creative Commons Attribution 4.0 International licence, as recorded in the arXiv licence metadata for this exact version and shown on the abstract page https://arxiv.org/abs/2408.11404v2. Full rights notice: licenses/arxiv-2408.11404-CC-BY-4.0.txt.\par\smallskip

\noindent\textit{Editorial scope.} Counted unit: the Proposition whose source label is \texttt{prop:unstable\_ii} in the article (source0.tex lines 1173--1179, source label \texttt{prop:unstable\_ii}), delivered together with the article's own complete original proof of it (source0.tex lines 1180--1233, one \texttt{proof} environment of 54 lines). No mathematics is written, repaired or re-derived: statement and proof are the article's own text line for line. Required context retained uncounted: lemma:lift\_section (lines 419--421); lemma:dimension\_section\_space (lines 537--543); lemma:aut\_vectorbundle (lines 574--588); lemma:difference\_loci (lines 1157--1162). Every definition, display and label of those lines is retained unchanged. Omitted from this package and not redistributed: 1--418, 422--536, 544--573, 589--1156, 1163--1172, 1234--1649. Nothing inside a retained excerpt is omitted or altered: every retained body is delivered exactly as the article's lines are written, so no omission receipt of any kind accompanies this package. Citations: the retained excerpts carry no citation command at all, so no bibliography entry of the article is transcribed and no reference is redistributed. Reference adaptation: none was needed and none was made; every reference inside the delivered unit resolves inside it, so no unresolved reference or citation is produced. Printed slips kept literal: no printed word, symbol, hypothesis or number of the counted statement or of its proof was altered. Layout and class: the article's own class is \texttt{extarticle} with packages including amsmath, amsthm, amssymb, color, graphicx, caption, mathtools, enumitem, verbatim, longtable, pifont, makecell, subcaption, xcolor; the wrapper is the lane's pinned \texttt{amsart} editorial wrapper at 10pt on A4, which restates the theorem-like environments the retained text needs and adds the article's own macro definitions that the retained text needs (\texttt{\textbackslash Aff}, \texttt{\textbackslash End}, \texttt{\textbackslash NN}, \texttt{\textbackslash OO}, \texttt{\textbackslash PP}, \texttt{\textbackslash cD}, \texttt{\textbackslash cH}, \texttt{\textbackslash cL}, \texttt{\textbackslash cV}, \texttt{\textbackslash op}, \texttt{\textbackslash tilde}), with extra wrapper packages (bm, bbm, dsfont, xspace, cancel, mathtools). The delivered numbering is the unit's own. Assets: no retained span contains a figure environment, a graphics-inclusion command, a PDF-page inclusion, a table or a TikZ picture; no image file is copied into this package, the package declares no asset dependency and no image file is redistributed with it. Licence: version arXiv:2408.11404v2 of this article is distributed under the Creative Commons Attribution 4.0 International licence, as recorded in the arXiv licence metadata for this exact version and shown on the abstract page https://arxiv.org/abs/2408.11404v2. A full rights notice is delivered at licenses/arxiv-2408.11404-CC-BY-4.0.txt. Characters: every retained excerpt is pure ASCII, so the delivered engine, XeLaTeX with its default Latin Modern text font, reports no missing character.
\begin{lemma}\label{lemma:lift_section}
Let $C$ be a smooth curve and $M, M' $ be line bundles on it. Let $s \in H^0(M')$ be a non-zero section with a reduced divisor $D$. Consider a non-trivial extension $\xi$ of $M'$ by $M$. Let $f: C \dashrightarrow \PP H^0(M^\vee \otimes M'\otimes K_C)$ be the map induced by the line bundle $M^\vee \otimes M'\otimes K_C$. The section $s$ lifts to a section of $V_\xi$ if and only if $[\xi] \in \PP H^0(M^\vee\otimes M' \otimes K_C)$ lies in the span of $f(D\setminus \op{Bs}(f))$. 
\end{lemma}

\begin{lemma}\label{lemma:dimension_section_space}
The dimension of $S_{\nu,\cL}$ is the following:
\begin{equation*}
\dim S_{\nu,\cL} = \op{max}_p \left(\dim Y_p + r_{\op{gen}}+p\right).
\end{equation*}
In particular if $\op{codim} Y_p \geq p $ for all $p \geq 1$ the dimension is $\dim Y + r_{\op{gen}}$. 
\end{lemma}

\begin{lemma}\label{lemma:aut_vectorbundle}
Let $(e,m)$ be a splitting type and $V=\bigoplus_{j=1}^l \OO_{\PP^1}(e_j)^{\oplus m_j}$. Then the automorphism group of $V$ consists of the following lower triangular block matrices:
\begin{equation*}
\begin{pmatrix}
A_{11} & 0 & \cdots & 0 \\
B_{21} & A_{22} & \cdots & 0 \\
\vdots & \vdots & \ddots &\vdots \\
B_{l1} & B_{l2} &\cdots& A_{ll}
\end{pmatrix}
\end{equation*}
Here the matrices $A_{jj}$ are elements of $GL_{m_j}(\Bbbk)$ and the $B_{ij}$ are arbitrary matrices with entries in $H^0(\OO_{\PP^1}(e_i - e_j)^{\oplus m_im_j})$. The dimension of $\op{Aut}(V)$ as an algebraic group is then
\begin{equation*}
\dim\op{Aut}(V) = \sum_{1\leq i \leq j \leq l}(e_j - e_i+1)m_im_j.
\end{equation*}
\end{lemma}

\begin{lemma}\label{lemma:difference_loci}
Let $0\leq d_2\leq d_1\leq g-1$ between $0$ and $g-1$. Let $\alpha$ be the difference map $C_{d_1}\times C_{d_2}\to \op{Pic}^{d_1-d_2}(C), (D_1,D_2)\mapsto \OO_C(D_1-D_2)$. For $r\geq 1$ the following holds: No component of $\alpha^{-1}(W^{r-1}_d(C))$ is entirely contained in $\alpha^{-1}(W^r_d(C))$. In particular we have
 \begin{equation*}
\op{codim} \alpha^{-1}(W^r_d(C)) \geq r+1.
\end{equation*} 
\end{lemma}

\begin{proposition}\label{prop:unstable_ii}
For $d_1 + d_2 < 1$ the locus $W_C^{0,0}(2,d_1,d_2)$ is empty. If $0 \leq d_2 < d_1 \leq g$ we have
\begin{equation*}
\dim W_C^{0,0}(2,d_1,d_2) = d_1 + 2d_2 -2 + 4(g-1)+1.
\end{equation*}
If $g+1\leq d_1$ we have a containment: $W_C^{0,0}(2,d_1,d_2) \subset W_C^{1,-1}(2,d_1,d_2)$.
\end{proposition}

\begin{proof}
Note first that effective line bundles have degree $\geq 0$ hence $d_2 \geq 0$ if $W^{(0,0)}_C(2,d_1,d_2)$ is non-empty. Furthermore we have $d_1 > d_2$. Hence $W^{(0,0)}_C(2,d_1,d_2)$ is empty if $d_1 + d_2 < 1$.

Now assume $0 \leq d_2< d_1 \leq g$.
We build a family of Higgs bundles over a scheme $M$ such that the induced map $M \to W_C^{0,0}(2,d_1,d_2)$ is surjective. We start with $M_1 = C_{d_1}\times C_{d_2}$. On $M_1 \times C$ we have two divisors $\cD_1,\cD_2$ which satisfy:
\begin{align*}
 \cD_{i|(D_1,D_2)\times C} = D_i
\end{align*}
for all $(D_1,D_2) \in M_1$ and $i =1,2$. We define $M_2$ to be $\Aff(\op{pr}_{1*}\OO(\cD_1)_{|\cD_2})$. Its dimension is $d_1 + 2d_2$. On $M_2 \times C$ we have an extension 
\begin{equation}\label{eq:univ_extension_ii}
0 \to \OO_{M_2\times C}(\cD_1) \to \cV \to \OO_{M_2\times C}(\cD_2) \to 0.
\end{equation}
 From Lemma \ref{lemma:lift_section} it follows that $h^0(\cV_{2|m\times C}) = 2$ for all $m \in M_2$. Put $M_3 = S_{\op{pr}_1,\End(\cV)\otimes \op{pr}_2^*(K_C)}$. We compute the dimension of $M_3$. We employ Lemma \ref{lemma:dimension_section_space}: To do this we need to find an upper bound for the dimension of the following loci:
\begin{equation*}
M_2^{(k)} := \left\{ m \in M_2 :  h^0(C,\op{End}\left(\cV_{1|m\times C}\right)\otimes K_C) = 4(g-1)+1+k\right\},\quad k\in \NN. 
\end{equation*}
By Riemann-Roch $m\in M_2$ is in $M_2^{(k)}$ if and only if $h^0(\End(\cV_{|m}))= k+1$. We decompose this locus in two parts:
\begin{align*}
M_{2,a}^{(k)} &= \left\{(D_1,D_2,\tilde{\xi}) \in M_2 : \begin{aligned}
&h^0(\OO(D_1-D_2)) = k \\
&\text{and the extension does not split}
\end{aligned} \right\},\\
M_{2,b}^{(k)} &= \left\{(D_1,D_2,\tilde{\xi}) \in M_2 : \begin{aligned}
&h^0(\OO(D_1-D_2)) = k-1 \\
&\text{and the extension splits}
\end{aligned} \right\}.
\end{align*}
One sees from Lemma \ref{lemma:difference_loci} that $M^{(k)}_{2,a}$ has codimension greater or equal to $k$. For $M^{(k)}_{2,b}$ we need to take the splitting into account. An element $\tilde{\xi}\in H^0(\OO(D_1)_{|D_2})$ induces the trivial extension if and only if it is in the kernel of $H^0(\OO(D_1)_{|D_2})\to H^1(\OO(D_1-D_2))$. Using a long exact sequence we see that the image has in general codimension $\deg D_2 -1$. 
We conclude 
\begin{align*}
\op{codim} M_{2,b}^{(k)} \geq k+d_2 - 2.
\end{align*}
If $d_2 \geq 2$ we have shown that $M_3$ has dimension $d_1 +2d_2 + 4(g-1)+1$. We will handle the cases $d_2 = 0$ and $d_2 = 1$ separately at the end of the proof. 

Let $M$ be the open subscheme of $M_3$ parametrizing stable Higgs bundles. We find the dimension of the image of $M\to \cH_C(2,d)$. Recall that an element of $M$ is a tuple
\begin{align*}
(D_1\in C_{d_1},D_2\in C_{d_2},\tilde{\xi} \in H^0(\OO(D_1)_{|D_2}),\phi).
\end{align*}
Here $\tilde{\xi}$ induces an extension with vector bundle $V$ and $\phi$ is a twisted endomorphism of $V$. Note that the morphism $H^0(\OO(D_1)_{|D_2}) \to H^1(\OO(D_1-D_2))$ has in the general case a one dimensional kernel $H^0(\OO(D_1))$. Furthermore we can multiply $\tilde{\xi}$ by a nonzero scalar and we obtain an isomorphic Higgs bundle. Therefore we have 
\begin{align}\label{eq:dimension_0,0unstable}
\dim W_C^{0,0}(2,d_1,d_2) &= \dim M - 2 \\
&= d_1 + 2d_2 -2 + 4(g-1)+1.
\end{align}

If $d_2 = 0$ all extensions parametrized by $M_2$ split because $D_2 = 0$ and therefore $H^0(\OO(D_1)_{|D_2})$ is zero too. Therefore $M^{(k)}_{2}$ equals $M^{(k)}_{2,b}$. This locus has codimension $\geq k-2$, where $k \geq 2$. We conclude that $M_3$ has dimension ${2d_2 +2  +4(g-1)+1}$. From lemma \ref{lemma:aut_vectorbundle} we see that for a general $m\in M_2$ we have $h^0(\End(\cV_{|m})) = 3$. This implies that the morphism $M \to \cH_C(2,d)$ generically has fibers of dimension $4$. Therefore $W_C^{0,0}(2,d_1,0)$ has dimension $d_1 -2  + 4(g-1)+1$. 

In case $d_2 = 1$ the extension generically splits. This can be seen as follows: Let ${(D_1,D_2,\tilde{\xi} \in H^0(\OO(D_1)_{|D_2}))}$ be a general point of $M_2$. This implies that $\OO(D_1-D_2)$ is not effective. Then we look at the long exact sequence of $0 \to \OO(D_1-D_2) \to \OO(D_1) \to \OO(D_1)_{|D_2} \to 0$:
\begin{equation*}
0 = H^0(\OO(D_1-D_2)) \to H^0(\OO(D_1)) \to H^0(\OO(D_1)_{|D_2}) \to H^1(\OO(D_1-D_2)).
\end{equation*} 
The map in the middle is an isomorphism and therefore the right map is $0$. Hence $\tilde{\xi}$ induces a split extension. One deduces that $M_3$ has dimension ${d_1 + 3 + 4(g-1) +1 }$. The group of automorphisms of $\OO(D_1)\oplus \OO(D_2)$ is in general two dimensional. Hence the fibers of $M\to \cH_C(2,d)$ are generically of dimension $3$. Therefore $W_C^{0,0}(2,d_1,d_2)$ has dimension $d_1 + 4(g-1) +1$ and $(V,\phi)$ is in $W_C^{1,-1}(2,d_1,d_2)$. 

For the last claim in the proof take a stable Higgs bundle $(V,\phi)$ with Harder-Narasimhan filtration $L_1 \subset V$ such that $\deg L_1 \geq g+1$. By the Riemann-Roch inequality we have $h^0(L_1)\geq 2$.
\end{proof}

\end{document}
